\chapter{Introduction}
\label{ch:intro}

This document is the \textbf{concise edition} of the B.Sc. project report. The
full report, with seven chapters, a glossary and a reproduction appendix, is a
separate file. Numbers, requirements and claim boundaries are the same in both
editions; what is shortened is explanation and repetition, not results.

Spoken conversation is the most natural form of human--machine interaction, but
until recently it was almost always built from a chain of independent
components. Models such as Moshi~\cite{defossez2024moshi} and
Freeze-Omni~\cite{wang2025freezeomni} showed that dialogue can be modelled as a
single ``speech-to-speech'' problem, with one model that listens and speaks at
the same time. This shift has happened for high-resource languages. For Persian,
none of the open models offers official support, and no conversational corpus
with interruption labels is publicly available.

The practical question of the project is: with the budget of a bachelor's
project, how far can one go in building a full-duplex Persian speech-to-speech
system, and what can be ``proven'' with auditable evidence?

\section{Problem Definition}
\label{sec:problem}

The defined goal is to design, implement and evaluate a prototype
speech-to-speech large language model for Persian. The system must be full
duplex: the user's microphone stays active while the reply plays, and the system
must be able to react to the user's ``interruption'' (barge-in). Three properties
separate this problem from conventional spoken question answering: listening and
speaking at the same time; a decision to stop playback within a few tens of
milliseconds; and continuity with what the user says at the moment of the
interruption, so that the user does not have to repeat the sentence.

\section{Motivation}

A speech interface can be the only effective way to use digital services for
users who are not comfortable with a Persian keyboard, for people with low
vision, and in particular for older adults. Persian speech recognition and
text-to-speech resources exist; a prominent example is
ManaTTS~\cite{qharabagh2025manatts}, with more than 86 hours of single-speaker
speech. These resources are read speech, however, and do not cover overlap,
interruptions or backchannels. On the other hand, current models such as
Qwen2.5-Omni~\cite{xu2025qwen25omni} need far more memory for their 7-billion
parameter version than the 12--24 GB range. Showing what is and is not
achievable within this range is itself an engineering result.

\section{Defined Requirements}
\label{sec:requirements}

The project definition sets seven specific requirements
(Table~\ref{tab:requirements}). The expected outputs are four items: a working
prototype, documented source code, a Persian conversational dataset, and
findings on interaction with older adults. For the fourth item, the measurement
protocol and infrastructure are ready; its empirical finding can only be
claimed once the human study is complete, so this report makes no claim about
the experience of older adults.

\begin{table}[ht]
\centering
\caption{The seven requirements of the project definition}
\label{tab:requirements}
\small
\begin{tabular}{clL{6.6cm}}
\toprule
\textbf{ID} & \textbf{Requirement} & \textbf{Specified measure} \\
\midrule
R1 & Persian speech-to-speech prototype & Persian spoken reply generated from
      Persian spoken input \\
R2 & Full-duplex operation & Listening during playback and reacting to
      interruptions \\
R3 & End-to-end latency & At most 500 ms from receiving speech to producing the
      spoken reply \\
R4 & Persian conversational dataset & 100--200 hours of speech with noise,
      overlap and interruptions \\
R5 & Interruption detection module & Accuracy above 80\% with energy, pitch and
      MFCC features \\
R6 & Evaluation on target hardware & A GPU with 12--24 GB of memory \\
R7 & Human evaluation & 5--10 Persian speakers, at least two aged over 60 \\
\bottomrule
\end{tabular}
\end{table}

\section{Two-Track Strategy and Challenges}
\label{sec:two-track}

Four challenges determined the final shape of the work: the lack of an official
training path for Persian in models such as Mini-Omni2~\cite{xie2024miniomni2}
and LLaMA-Omni2~\cite{fang2025llamaomni2}; the lack of a labelled Persian
conversational corpus; the gap between reducing offline loss and valid
autoregressive generation; and the boundary between engineering evidence and
scientific claims.

To reduce risk, two tracks were pursued in parallel. The direct track is
low-rank adaptation of Moshika with the official trainer of the Moshi
family~\cite{moshifinetune2024}. The cascade track is a chain of Persian speech
recognition, a text language model and text-to-speech, on top of which
full-duplex control and interruption logic are built separately. Both tracks
used the same dataset. The outcome was asymmetric: the cascade track reached a
streaming full-duplex system that meets the numeric latency and interruption
criteria in automated tests, and the direct track is reported as a negative
result together with its measured cause.

\section{Contributions}
\label{sec:contributions}

\begin{enumerate}
\item \textbf{A streaming full-duplex Persian voice assistant.} Open
      microphone, speculative end-of-speech detection and clause-by-clause
      playback; in browser loopback, a mean latency of 372 ms and a 90th
      percentile of 431 ms with 11.1~GiB of GPU memory.

\item \textbf{Two-stage interruption detection.} Ducking after one frame and a
      lexical confirmation with echo rejection; accuracy 0.933/F1 0.880 in
      simulation and 0.825 in the browser, against 0.574 for an acoustic
      classifier alone.

\item \textbf{An audited dataset.} 6,754 pairs with 108.584 hours of audio,
      session-based splits and zero group leakage.

\item \textbf{Evidence-based component selection.} A final semantic test
      (3.150 versus 1.425) and re-selection of the recognizer, responder and
      synthesizer on fresh panels.

\item \textbf{A measured cause for the failure of the direct track.} Seven
      generations of Moshika and one speech-input model; the tokenizer
      bottleneck is real but not the only obstacle.

\item \textbf{An evidence-first method.} Protocols were registered before
      running, test splits were opened once, and artifacts are write-once.
\end{enumerate}

Latency and interruption evidence is automated and on an H100, and the results
of the listening study are not yet available.

\section{Structure of This Report}

Chapter~\ref{ch:background} condenses the necessary concepts and the position of
the work. Chapter~\ref{ch:method} describes how the data were built, the cascade
architecture, the interruption detector, the direct-track experiments and the
evidence classes. Chapter~\ref{ch:results} reports the measured results with the
claim boundary of each. Chapter~\ref{ch:conclusions} summarizes the
interpretation, limitations and future work.
